% paper/sections/08f_ppe_bc.tex
%
% §8.4  境界条件と数値安定性
%
\subsection{境界条件と数値安定性}
\label{sec:ppe_bc_stability}
\begin{itemize}
    \item \textbf{固体壁・流入出：} $\partial p/\partial n = 0$（ノイマン条件）．CCD ブロック方程式への組み込み手順（境界行を $p'_0=0$ で置き換え，変密度項の消滅）は §\ref{sec:ccd_neumann_bc} を参照．
    \item \textbf{自由表面：} $p = p_{\mathrm{atm}}$（ディリクレ条件）
    \item \textbf{全周ノイマンの場合：} 解の一意性を保つため，1点の圧力を固定するか平均値を0にする拘束を加える（§\ref{sec:pinpoint}参照）．
    \item \textbf{周期境界：} 全辺を周期とする場合は FVM 行列に折り返しフェイスおよびゴーストノード拘束を組み込む（§\ref{sec:fvm_periodic} 参照）．零空間対策は壁面 BC と同一（1点ピン条件）．
\end{itemize}

\label{box:neumann_unit_test}
Neumann 境界条件 $\partial(\delta p)/\partial n = 0$ を CCD スキームに正しく組み込んだか確認するための最小検証セット：

\begin{enumerate}
  \item \textbf{無負荷試験（Null test）：}
        $q_h \equiv 0$（右辺ゼロ），初期値 $(\delta p)^0 = 0$ として反復を実行し，
        $\|\delta p^m\|_2 \to 0$（機械精度）が達成されることを確認．
        ゼロでなければ境界ブロックの係数設定に誤りあり．
  \item \textbf{既知解収束テスト（Manufactured Solution）：}
        Neumann BC に整合する既知解，例えば
        $p_\text{exact}(x,y) = \cos(\pi x)\cos(\pi y)$
        （$\partial p/\partial n = 0$ を4辺全てで満たす）に対し，
        $q_h = \mathcal{L}_h p_\text{exact}$ を右辺として与え，
        収束後の誤差が $\|p^m - p_\text{exact}\|_2 = \Ord{h^6}$（格子収束）であることを確認．
        $\Ord{h^2}$ のみなら境界誤差が伝播しており，§\ref{sec:ccd_bc} の $L^2$ ノルム重み付け解析（式~\eqref{eq:bc_L2_weighting}）を参照．
  \item \textbf{零モード（定数解）のフィルタリング確認：}
        $q_h \equiv 0$ で $(\delta p)^0 = c$（定数）を初期値として反復し，
        ピン条件（§\ref{sec:pinpoint}）が正しく機能して $p_{i_0,j_0} \to 0$ が保たれることを確認．
        定数モードが収束を妨げる場合は零モード除去の実装を確認すること．
\end{enumerate}

\subsubsection{FVM PPE 行列の周期境界条件対応}
\label{sec:fvm_periodic}

壁面（ノイマン）BC に対する FVM 離散化（式~\eqref{eq:af_exact}）は，
周期境界を持つ計算領域ではそのまま適用できない．
本稿の実装では以下の3点を変更して周期 FVM 行列を組み立てる．

\paragraph{（1）折り返しフェイスの追加}
通常の軸 $x$ 方向内部フェイスはノード $i\,(i=0,\ldots,N-2)$ とノード $i+1$ を結ぶ
$N-1$ 本で構成される．
周期 BC ではさらに\textbf{折り返しフェイス（wrap face）}
としてノード $N_x-1$ とノード $0$ を結ぶフェイスを追加する：
\begin{equation}
  a_{N_x - 1/2} = \frac{2}{\rho_{N_x-1} + \rho_0}
  \label{eq:af_wrap}
\end{equation}
このフェイスにより，スパース行列は完全な循環ステンシルを持つ．

\paragraph{（2）境界半格子体積補正の省略}
壁面 BC では境界ノード（$i=0$ または $i=N$）の制御体積が
$h/2$（半格子幅）であるため，フェイス係数の寄与を2倍に補正する．
周期 BC ではすべてのノードが同じ格子幅 $h$ の制御体積を持つため，
この補正は行わない（$a_f / h^2$ を一様に適用）．

\paragraph{（3）ゴーストノードの同一性拘束}
本稿の格子は各軸方向に $N+1$ ノード（インデックス $0,1,\ldots,N$）を持つ．
インデックス $N$ のノードは物理的に $0$ の周期像（ゴーストノード）であり，
独立な自由度ではない．
これらのゴーストノードに対応するスパース行列の行を
\textbf{同一性拘束行（identity-minus）}で置き換える：
\begin{equation}
  p_{\mathrm{ghost}} - p_{\mathrm{src}} = 0
  \quad\Longleftrightarrow\quad
  A[\mathrm{ghost},\,\mathrm{ghost}] = 1,\quad
  A[\mathrm{ghost},\,\mathrm{src}]   = -1,\quad
  b[\mathrm{ghost}] = 0
  \label{eq:ghost_identity}
\end{equation}
ここで $\mathrm{src}$ はゴーストノードに対応する物理ノード（インデックス $0$）を指す．
各軸方向に独立にゴーストノードを列挙し，複数軸で同時にゴーストとなるコーナー節点は
最初の軸の拘束のみを採用する（重複適用を避けるため）．

\begin{tcolorbox}[warnbox, title={周期 FVM 行列と前処理：ILU の充填不足に注意}]
  式~\eqref{eq:ghost_identity} の同一性拘束行と FVM 行が混在した疎行列に対して，
  充填率 \texttt{fill\_factor}$=1$ の不完全 LU（ILU(0)）を前処理として使用すると，
  BiCGSTAB の収束が著しく遅延する（FVM 行は対角優位だが同一性拘束行は異なる構造を持つため）．
  実装では周期 BC に対して \texttt{fill\_factor}$=10$ に引き上げることで安定収束を確認した．
  直接法（SuperLU / spsolve）はこの問題を回避できるため，
  収束診断の際は \texttt{ppe\_solver\_type='lu'} で動作確認を行うと有効である．
\end{tcolorbox}

\subsubsection{高密度比における条件数と数値安定性}
\label{sec:ppe_condition_number}

気液界面を挟む密度比 $\rho_l/\rho_g \approx 10^3$（水/空気相当）は
PPE 演算子の条件数に直接影響する．
可変係数演算子 $\mathcal{L}_h$ の係数 $a_{i+1/2} = 2/(\rho_i+\rho_{i+1})$ は，
液相（$\rho \approx \rho_l$）と気相（$\rho \approx \rho_g$）の格子フェイスで
$O(\rho_l/\rho_g)$ の比率を持つため，行列の条件数は
\begin{equation}
  \kappa(\mathcal{L}_h) = O\!\left(\frac{\rho_l}{\rho_g}\cdot\frac{1}{h^2}\right)
  \approx O\!\left(\frac{10^3}{h^2}\right)
  \label{eq:ppe_cond_number}
\end{equation}
と見積もられる．

\noindent\textbf{安定化戦略：}
\begin{itemize}
  \item \textbf{調和平均フェイス係数（式~\eqref{eq:af_exact}）：}
        単純算術平均（$(\rho_L+\rho_R)/2$ の逆数）と比較して，
        調和平均 $2/(\rho_L+\rho_R)$ は界面フラックスに物理的根拠を持ち，
        条件数の過大評価を防ぐ（付録~\ref{app:fvm_face_coeff}）．
  \item \textbf{GFM によるジャンプの明示的処理（§\ref{sec:gfm}）：}
        CSF 体積力を用いず Young--Laplace ジャンプ $[p]_\Gamma = \kappa/We$ を
        右辺補正として組み込むことで，圧力場が界面で鋭利に保たれ，
        反復ソルバーの残差フロアへの影響が最小化される．
  \item \textbf{ゲージ固定点の対称性（KL-11；§\ref{sec:pinpoint}）：}
        ピン固定点を中心ノード $(N/2, N/2)$ に設定することで（コーナーピン禁止），
        高密度比条件下でもグローバル非対称性 $O(\|\mathrm{rhs}\|)$ を回避する
        （ASM-008 対称性修正，2026-03-22 適用済み）．
  \item \textbf{直接 LU ソルバーへのフォールバック：}
        LGMRES が収束しない場合（ASM-005），\texttt{spsolve}（直接 LU）で
        条件数の影響を受けず解を得られる（実装は付録~\ref{sec:ppe_pseudotime}）．
\end{itemize}
